\chapter{Quantum impurities and nanostructures}
\setwlchapter{impurities}
\section{A local level embedded in a band}
For one spinful impurity the Anderson Hamiltonian reads
\[
 H=\sum_{k\sigma}\epsilon_k n_{k\sigma}+\epsilon_d n_d+U D_d+
 \sum_{k\sigma}(V_k c^\dagger_{k\sigma}d_\sigma+\hc).
\]
At $V_k=0$ the local moment exists for $\epsilon_d<0<\epsilon_d+U$.
A bath adds infinitely many low-energy configurations; it cannot in
general be represented by a constant shift of the impurity level.
Already at $U=0$ the retarded hybridization is
$\Delta(z)=\sum_k|V_k|^2/(z-\epsilon_k)$, whose imaginary part
describes the widths of levels only in the continuum limit.
\casestudy{leads}{Hybridization to a local lattice orbital}
\casestudy{resonance}{Frequency-dependent hybridization with one bath level}

\section{From valence fluctuations to Kondo exchange}
For initial and final bath states close to the Fermi level, eliminate
empty and doubly occupied impurity states in second order. For real
constant $V$ the exchange multiplying $\vS_d\cdot\mathbf s(0)$ is
$J=2V^2[1/(-\epsilon_d)+1/(\epsilon_d+U)]$. Both virtual paths favour
the impurity--bath singlet. Their \emph{difference} controls potential
scattering, which vanishes at $\epsilon_d=-U/2$. The spin density here
uses Pauli matrices divided by two; changing that definition changes
the quoted $J$. This elimination is controlled when both charge
excitation energies are large on the hybridization scale. It does not
directly compute the nonperturbative Kondo scale.
\casestudy{anderson}{Checking exchange against exact diagonalization}

\section{Wilson chain and NRG Ljubljana}
A uniform finite bath has a lowest level spacing that competes with
the screening scale. NRG instead divides the band into logarithmic
intervals and transforms the retained combinations into a chain.
Only the first chain orbital couples to the impurity. Far along a
metallic Wilson chain the hopping typically decays as
$\Lambda^{-n/2}$ for logarithmic discretization parameter $\Lambda>1$.
Successive diagonalizations can then discard high-energy states relative
to the next hopping. Check a discretization and truncation sequence;
a two-site singlet in an untruncated test is a normalization check,
not evidence by itself of continuum screening.

In the public \href{https://github.com/rokzitko/nrgljubljana}{NRG Ljubljana}
tree, \path{test/c++/test42_model_1orb/model.m} and its companion
\texttt{param} give a complete one-channel initializer example.
\fn{def1ch[1]} sets up the impurity and bath end; the local term uses
\fn{number[d[]]} and \fn{hubbard[d[]]}, the contact uses
\fn{hop[f[0],d[]]}, and the host term is included once. The initializer
derives the contact normalization from the discretized bath; inserting
the continuum width $\Gamma$ as the hopping would have the wrong
normalization. The setting \texttt{symtype=QS} assumes charge U(1) and
spin SU(2) invariance; choose different
quantum numbers if the model has pairing, spin anisotropy or a magnetic
field. Specify measured observables too: an energy spectrum contains
no information about matrix elements of a local spin flip.

\section{Charge stability and superconducting islands}
In a double dot the electrostatic term $U_{12}n_1 n_2$ competes
with the on-site $U$; near the $(1,1)$ configuration the virtual
charge-transfer cost is $U-U_{12}$ for equivalent sites. The same
second-order hopping that mixes $(2,0)$ and $(0,2)$ with a singlet
therefore controls the singlet--triplet splitting. Close to a charge
degeneracy both charge configurations must be retained.
For multi-orbital dots a high-spin atomic multiplet can replace the
spin-$1/2$ moment (see Chapter~10 and Ref.~\cite{horvat}).

A superconducting island brings both pairing and charging energy
$E_C(N-N_g)^2$. A static pair field can model proximity to a bulk
reservoir but is not a number-conserving description of a floating
island. Compare even and odd \emph{total} parity sectors, include
the energy of a residual quasiparticle, and distinguish a spin's
local response from its total-spin quantum number. The latter
distinction enters the impurity Knight shift~\cite{pavesic-knight};
crystal fields provide another way of altering the impurity response
without changing its charge~\cite{zitko-anisotropic}.
\begin{exercise}[Where does the pole weight go?]\label{ex:resonance}
Hybridize a noninteracting level at zero energy with one bath level
also at zero energy by a real $V>0$. Find both poles of the impurity
propagator, their residues per spin, and the limit as $V\to0$.
\end{exercise}
